% Generated draft: requires mathematical and source review.
\documentclass[11pt]{article}
\usepackage[margin=25mm]{geometry}
\usepackage{fontspec}
\setmainfont{Latin Modern Roman}
\usepackage{amsmath,amssymb,mathtools,unicode-math}
\setmathfont{Latin Modern Math}
\usepackage{enumitem,needspace,xurl,hyperref}
\setlength{\parindent}{0pt}
\setlength{\parskip}{5pt}
\setlength{\emergencystretch}{3em}
\newcommand{\N}{\mathbb{N}}
\newcommand{\Z}{\mathbb{Z}}
\newcommand{\Q}{\mathbb{Q}}
\newcommand{\R}{\mathbb{R}}
\newcommand{\C}{\mathbb{C}}
\newcommand{\F}{\mathbb{F}}
\newcommand{\A}{\mathbb{A}}
\newcommand{\PP}{\mathbb{P}}
\newcommand{\E}{\mathbb{E}}
\newcommand{\eps}{\varepsilon}
\newcommand{\dd}{\,\mathrm{d}}
\begin{document}
\section*{UnsolvedMath ID: 700632}
\textbf{Source problem code:} catalogue entry
\subsection{Definitions and mathematical statement}
A finite group \(G\) is quasisimple if \(G=[G,G]\) and \(G/Z(G)\) is a nonabelian simple group, where \([G,G]\) is the commutator subgroup and \(Z(G)\) is the centre. Let \(p\) be a prime and let \(k\) be a field of characteristic \(p\). A finite-dimensional \(p\)-modular representation over \(k\) is a homomorphism \(\rho\colon G\to\operatorname{GL}_k(V)\), where \(V\) is a nonzero finite-dimensional \(k\)-vector space. Its degree is \(\dim_k V\). It is faithful if \(\ker\rho=\{1\}\), and absolutely irreducible if \(V\otimes_k\overline{k}\) is irreducible as a representation of \(G\) over an algebraic closure \(\overline{k}\) of \(k\). Two representations over \(k\) are equivalent if there is an invertible \(k\)-linear map intertwining their actions.

Write \(\mathcal R_k(G)\) for the equivalence classes of faithful absolutely irreducible representations of \(G\) over \(k\). When \(\mathcal R_k(G)\ne\varnothing\), define
\[
 d_k(G)=\min\{\dim_k V:[\rho\colon G\to\operatorname{GL}_k(V)]\in\mathcal R_k(G)\}
\]
and
\[
 \mathcal M_k(G)=\{[\rho\colon G\to\operatorname{GL}_k(V)]\in\mathcal R_k(G):\dim_k V=d_k(G)\}.
\]
\textbf{Problem (A. S. Kondratiev).} For every \emph{(known)} finite quasisimple group \(G\) and every prime \(p\), find the faithful \(p\)-modular absolutely irreducible linear representations of minimal degree. With a coefficient field \(k\) fixed, this means determining \(d_k(G)\) and \(\mathcal M_k(G)\) whenever \(\mathcal R_k(G)\) is nonempty.

The coefficient field is retained as an explicit parameter: the source does not specify it or whether equivalence is intended over an algebraic closure. The source also does not enumerate the groups covered by ``(known)'' or specify the date governing that qualification. If \(\mathcal R_k(G)\) is empty, the displayed minimum is undefined; the source does not explicitly discuss this case. These qualifications remain unresolved.
\subsection{Short English statement}
For every finite quasisimple group covered by the source's qualification ``(known)'', and every prime \(p\), determine the faithful absolutely irreducible representations in characteristic \(p\) having the smallest possible dimension. The coefficient-field convention, the scope of ``known'', and the treatment of cases with no such representation require clarification.
\subsection{Sources}
\begin{itemize}
\item Local UnsolvedMath registry and supplied complete source record: permanent ID \(700632\), exact source problem code \texttt{catalogue entry}, title \emph{Kourovka Notebook Problem 10.15}. The record is labelled \texttt{open}; this is a supplied status label, not an independently verified mathematical or literature assessment.
\item \emph{The Kourovka Notebook (Unsolved Problems in Group Theory)}, original problem 10.15, 10th issue (1986), proposed by A. S. Kondratiev. The supplied record identifies the 21st-issue edition (2026), \texttt{arXiv:1401.0300v48}.
\item The supplied collection description attributes editing to E. I. Khukhro and V. D. Mazurov and describes coverage of Issues 1--21 (1965--2026). The references and edition information above are supplied attribution; they have not been independently consulted or verified for this draft.
\end{itemize}
\end{document}
